\chapter{Intégrales curvilignes et intégrales multiples}\label{ch:b2:multint}

Ce chapitre étend l'intégration des segments aux courbes et aux
domaines du plan et de l'espace. Les \omterm{def:b2:multint:lineint}{intégrales curvilignes} intègrent
une \emph{\omterm{def:b2:multint:lineint}{forme différentielle}} $P\,\dd x + Q\,\dd y$ le long d'un arc
orienté~; les intégrales doubles et triples intègrent des fonctions sur
des domaines de dimension deux et trois. Les deux théories se
rencontrent dans le \emph{théorème de Green-Riemann}, le théorème
fondamental de l'analyse en dimension deux, et l'outil de calcul
principal tout au long est la \emph{formule de changement de
variables}, dont le facteur de distorsion est la valeur absolue du
\omterm{def:b2:linalg:det}{déterminant} jacobien.

\section{Intégrales curvilignes}

\begin{definition}[Forme différentielle~; intégrale curviligne]\label{def:b2:multint:lineint}
Soit $U \subseteq \R^2$ un \omterm{def:b2:metric:topology}{ouvert}. Une \emph{forme
différentielle}\index{forme différentielle} de
degré $1$ et de classe $\mathcal{C}^0$ sur $U$ est une expression
$\omega = P\,\dd x + Q\,\dd y$ avec $P, Q \colon U \to \R$
\omterm{def:b2:metric:continuity}{continues} --- formellement, une application \omterm{def:b2:metric:continuity}{continue} de $U$ dans le
\omterm{def:b2:linalg:dual}{dual} de $\R^2$, $\omega(M) = P(M)\,e_1^* + Q(M)\,e_2^*$. Pour un
arc $\mathcal{C}^1$ $\gamma \colon [a, b] \to U$, $\gamma(t) =
(x(t), y(t))$, l'\emph{intégrale curviligne}\index{intégrale curviligne} de
$\omega$ le long de $\gamma$
est
\[
\int_\gamma \omega
= \int_a^b \Bigl(P(\gamma(t))\,x'(t)
+ Q(\gamma(t))\,y'(t)\Bigr)\,\dd t .
\]
Les définitions s'étendent mot pour mot à $\R^3$ (formes $P\,\dd x +
Q\,\dd y + R\,\dd z$) et aux arcs $\mathcal{C}^1$ par morceaux (somme
sur les morceaux).
\end{definition}

\begin{proposition}[Invariance et orientation]\label{prop:b2:multint:lineinv}
L'\omterm{def:b2:multint:lineint}{intégrale curviligne} est inchangée par un \omterm{def:b2:curves:reparam}{changement de paramètre}
$\mathcal{C}^1$ croissant, et change de signe sous un changement
décroissant. Elle ne dépend donc que de l'\omterm{def:b2:curves:reparam}{arc géométrique}
\emph{orienté}.
\end{proposition}

\begin{proof}
Si $\theta \colon [c, d] \to [a, b]$ est un \omterm{def:b2:curves:reparam}{changement de paramètre} et
$\tilde\gamma = \gamma \circ \theta$, alors par la règle de la chaîne
et le changement de variables à une variable $t = \theta(u)$,
\[
\int_{\tilde\gamma}\omega
= \int_c^d \bigl(P(\gamma(\theta(u)))\,x'(\theta(u))
+ Q(\gamma(\theta(u)))\,y'(\theta(u))\bigr)\,\theta'(u)\,\dd u
= \pm\int_a^b \bigl(Px' + Qy'\bigr)(t)\,\dd t ,
\]
avec le signe $+$ si $\theta$ est croissant ($\theta(c) = a$) et $-$
s'il est décroissant (les bornes s'échangent).
\end{proof}

\begin{example}[Travail d'une force~; circulation]\label{ex:b2:multint:work}
Si $F = (P, Q)$ est un champ de forces, $\int_\gamma P\dd x + Q\dd y =
\int_a^b \langle F(\gamma(t)), \gamma'(t)\rangle\,\dd t$ est le
\emph{travail} de $F$ le long de $\gamma$. Pour $\omega = -y\,\dd x +
x\,\dd y$ le long du cercle unité parcouru dans le sens direct
$\gamma(t) = (\cos t, \sin t)$~:
\[
\int_\gamma \omega
= \int_0^{2\pi}\bigl((-\sin t)(-\sin t)
+ \cos t\cos t\bigr)\,\dd t = 2\pi ,
\]
deux fois l'\omterm{def:b2:multint:domain}{aire} délimitée --- un premier indice de Green-Riemann.
\end{example}

\begin{example}[Une intégrale, deux paramétrages, un
piège de signe]\label{ex:b2:multint:semicircle}
Calculons $\int_\gamma x\,\dd y$ le long du demi-cercle unité
supérieur de $(1, 0)$ à $(-1, 0)$. Avec $\gamma(t) =
(\cos t, \sin t)$, $t \in \intcc0\pi$~:
\[
\int_0^\pi\cos t\cdot\cos t\,\dd t = \frac\pi2 .
\]
Avec le paramétrage par le graphe $x \mapsto (x, \sqrt{1 -
x^2})$, $x$ allant de $1$ à $-1$ (attention au sens~!)~:
\[
\int_1^{-1}x\cdot\frac{-x}{\sqrt{1 - x^2}}\,\dd x
= \int_{-1}^{1}\frac{x^2}{\sqrt{1 - x^2}}\,\dd x
= \frac\pi2
\]
($x = \sin u$ le ramène à une intégrale de Wallis). Même valeur,
comme le garantit la~\cref{prop:b2:multint:lineinv} --- mais seulement
parce que les deux parcours vont de $(1,0)$ à $(-1,0)$~; inverser le
sens de parcours change le signe. Refermer le chemin le long de l'axe
des $x$ (où $\dd y = 0$) n'ajoute rien, et le total $\frac\pi2$
est l'\omterm{def:b2:multint:domain}{aire} du demi-disque~: la première occurrence des formules
aire-bord de Green-Riemann ci-dessous.
\end{example}

\begin{definition}[Formes exactes et fermées]\label{def:b2:multint:exact}
La forme $\omega = P\,\dd x + Q\,\dd y$ de classe $\mathcal{C}^0$
est \emph{exacte}\index{forme exacte} sur $U$ s'il existe $f \in \mathcal{C}^1(U)$
(un \emph{potentiel}\index{potentiel}) tel que $\omega = \dd f$, c.-à-d.\ $P = f_x$ et $Q
= f_y$. Une forme $\mathcal{C}^1$ est \emph{fermée}\index{forme fermée} si $P_y = Q_x$ sur
$U$.
\end{definition}

\begin{theorem}[Théorème fondamental pour les intégrales
curvilignes]\label{thm:b2:multint:ftc}
Si $\omega = \dd f$ est \omterm{def:b2:multint:exact}{exacte} et si $\gamma$ est un arc $\mathcal{C}^1$
par morceaux dans $U$ de $A$ à $B$, alors
\[
\int_\gamma \omega = f(B) - f(A) .
\]
En particulier, l'intégrale d'une \omterm{def:b2:multint:exact}{forme exacte} le long de tout arc
fermé est nulle, et toute forme $\mathcal{C}^1$ \omterm{def:b2:multint:exact}{exacte} est \omterm{def:b2:multint:exact}{fermée}.
\end{theorem}

\begin{proof}
$\frac{\dd}{\dd t}f(\gamma(t)) = f_x(\gamma(t))x'(t) +
f_y(\gamma(t))y'(t)$ par la règle de la chaîne
(\cref{ch:b2:diffcalc}), donc l'intégrande de la~%
\cref{def:b2:multint:lineint} est la dérivée de $t \mapsto
f(\gamma(t))$, et le théorème fondamental de l'analyse donne le
résultat sur chaque morceau~; les valeurs intermédiaires se
télescopent. Le fait que les formes $\mathcal{C}^1$ \omterm{def:b2:multint:exact}{exactes} soient
\omterm{def:b2:multint:exact}{fermées} est le théorème de Schwarz~: $P_y =
f_{xy} = f_{yx} = Q_x$.
\end{proof}

\begin{example}[Reconstruction d'un potentiel]\label{ex:b2:multint:potential}
Soit $\omega = y\,\eu^{xy}\,\dd x + (x\,\eu^{xy} + 2y)\,\dd y$
sur $\R^2$. Elle est \omterm{def:b2:multint:exact}{fermée}~: les deux dérivées croisées valent
$\eu^{xy}(1 + xy)$. Pour trouver un \omterm{def:b2:multint:exact}{potentiel}, intégrons $P$ en
$x$ à $y$ fixé~:
\[
f(x, y) = \int y\,\eu^{xy}\,\dd x = \eu^{xy} + c(y),
\]
puis ajustons $c$ en identifiant $f_y$~: $x\,\eu^{xy} + c'(y) =
x\,\eu^{xy} + 2y$ donne $c(y) = y^2$. Donc $f(x,y) = \eu^{xy} +
y^2$, et pour tout arc $\mathcal C^1$ par morceaux de $(0,0)$
à $(1,1)$,
\[
\int_\gamma\omega = f(1,1) - f(0,0) = (\eu + 1) - 1 = \eu ,
\]
indépendamment du chemin --- la recette en deux temps (intégrer
en $x$, corriger en $y$) est la réciproque pratique du~%
\cref{thm:b2:multint:ftc} sur les domaines où les \omterm{def:b2:multint:exact}{formes fermées} sont
\omterm{def:b2:multint:exact}{exactes}.
\end{example}

\begin{example}[Fermée n'implique pas exacte]\label{ex:b2:multint:angleform}
Sur $U = \R^2 \setminus \{0\}$, la \emph{forme angulaire}
\[
\omega = \frac{-y\,\dd x + x\,\dd y}{x^2 + y^2}
\]
est \omterm{def:b2:multint:exact}{fermée} (calcul direct~: $P_y$ et $Q_x$ valent tous deux
$\frac{y^2 - x^2}{(x^2+y^2)^2}$), mais son intégrale le long du cercle
unité vaut $2\pi \neq 0$ (même calcul qu'à l'\cref{ex:b2:multint:work},
divisé par $1$)~: $\omega$ n'est pas \omterm{def:b2:multint:exact}{exacte}
sur $U$. Localement, $\omega = \dd\theta$ pour une détermination
$\theta$ de l'angle polaire~; l'obstruction est globale --- l'angle
ne peut pas être défini continûment autour du point retiré. Sur les
domaines sans trous, la pathologie disparaît~: sur un \omterm{def:b2:metric:topology}{ouvert}
\emph{étoilé}, toute forme $\mathcal{C}^1$ \omterm{def:b2:multint:exact}{fermée} est \omterm{def:b2:multint:exact}{exacte}
(lemme de Poincaré, \cref{exo:b2:multint:8}).
\end{example}

\section{Intégrales doubles}

Nous admettons la théorie de l'intégrale de Riemann à une
variable (volume de Licence~1, et \cref{ch:b2:integration}) et
esquissons sa version à deux variables. Une fonction $f$ \omterm{def:b2:metric:continuity}{continue} sur
un rectangle $R = [a, b] \times [c, d]$ a une intégrale double
$\iint_R f$, définie par des sommes de Riemann sur des grilles
exactement comme à une variable, et calculée par itération~:

\begin{theorem}[Fubini sur un rectangle]\label{thm:b2:multint:fubini}
Pour $f$ \omterm{def:b2:metric:continuity}{continue} sur $R = [a,b] \times [c,d]$,
\[
\iint_R f
= \int_a^b \Bigl(\int_c^d f(x, y)\,\dd y\Bigr)\dd x
= \int_c^d \Bigl(\int_a^b f(x, y)\,\dd x\Bigr)\dd y .
\]
\end{theorem}

\begin{proof}
Posons $F(x) = \int_c^d f(x, y)\,\dd y$. La continuité uniforme de $f$
sur le \omterm{def:b2:metric:compact}{compact} $R$ rend $F$ \omterm{def:b2:metric:continuity}{continue} (majoration par domination~:
$\abs{F(x) - F(x')} \leq (d - c)\sup_y\abs{f(x,y) - f(x',y)}$).
Subdivisons maintenant $[a,b]$ et $[c,d]$ en $n$ parts égales, ce qui
donne une grille de cellules $R_{ij}$ d'\omterm{def:b2:multint:domain}{aire} $\Delta x\,\Delta y$. Sur
chaque cellule,
$\inf_{R_{ij}} f \cdot \Delta x \Delta y \leq
\int_{x_{i-1}}^{x_i}\int_{y_{j-1}}^{y_j} f(x,y)\,\dd y\,\dd x \leq
\sup_{R_{ij}} f \cdot \Delta x \Delta y$ par croissance de l'intégrale
à une variable (appliquée deux fois). En sommant sur les cellules,
l'intégrale itérée $\int_a^b F$ est encadrée entre les sommes de
Riemann inférieure et supérieure de la grille~; par continuité
uniforme, les deux sommes convergent vers la valeur commune qui
définit $\iint_R f$ lorsque $n \to
\infty$. Le même argument s'applique en échangeant les rôles de $x$
et $y$, donc les deux intégrales itérées sont égales à $\iint_R f$.
\end{proof}

\begin{remark}
La continuité sur un rectangle \omterm{def:b2:metric:compact}{compact} joue un rôle essentiel dans
la preuve de Fubini~: c'est elle qui fournit la continuité uniforme
encadrant les sommes de Riemann. Pour des intégrandes plus
sauvages, l'énoncé est réellement en défaut --- il existe des
fonctions dont les deux intégrales itérées existent et diffèrent. Le
théorème général honnête, avec l'intégrabilité pour seule hypothèse,
est le théorème de Fubini pour l'intégrale de Lebesgue, prouvé dans le
volume de Licence~3~; tout ce chapitre reste dans le cadre \omterm{def:b2:metric:continuity}{continu} où
la preuve élémentaire ci-dessus est complète.
\end{remark}

\begin{definition}[Domaines élémentaires]\label{def:b2:multint:domain}
Un domaine $D \subseteq \R^2$ est \emph{élémentaire selon $y$} si
\[
D = \{(x, y) : a \leq x \leq b,\
\varphi_1(x) \leq y \leq \varphi_2(x)\}
\]
avec $\varphi_1 \leq \varphi_2$ \omterm{def:b2:metric:continuity}{continues} sur $[a,b]$
(élémentaire selon $x$~: symétriquement). Pour $f$ \omterm{def:b2:metric:continuity}{continue} sur un
domaine $D$ élémentaire selon $y$,
\[
\iint_D f
= \int_a^b\Bigl(
\int_{\varphi_1(x)}^{\varphi_2(x)} f(x,y)\,\dd y\Bigr)\dd x ,
\]
et l'on vérifie (en prolongeant $f$ par un argument d'approximation, ou
en subdivisant) que lorsque $D$ est élémentaire dans les deux
directions les deux intégrales itérées coïncident. Les domaines
découpés en un nombre fini de morceaux élémentaires se traitent par
additivité, et l'\emph{aire}\index{aire!d'un domaine plan} de $D$ est $\operatorname
{Area}(D) = \iint_D 1$.
\end{definition}

\begin{example}\label{ex:b2:multint:triangle}
Sur le triangle $D = \{0 \leq x \leq 1,\ 0 \leq y \leq x\}$~:
\[
\iint_D xy \,\dd x\,\dd y
= \int_0^1 x\Bigl(\int_0^x y\,\dd y\Bigr)\dd x
= \int_0^1 x\cdot\frac{x^2}{2}\,\dd x = \frac18 .
\]
En échangeant l'\omterm{def:b2:structures:generated}{ordre} ($x$ allant de $y$ à $1$)~: $\int_0^1
y\bigl(\int_y^1 x\,\dd x\bigr)\dd y = \int_0^1
y\,\frac{1 - y^2}{2}\,\dd y = \frac18$ --- même valeur, calcul
différent~: bien choisir l'\omterm{def:b2:structures:generated}{ordre} d'intégration, c'est la moitié du
métier.
\end{example}

\begin{example}[Quand un seul ordre fonctionne]\label{ex:b2:multint:swaporder}
Calculons $I = \displaystyle\int_0^1\!\!\int_x^1
\eu^{y^2}\,\dd y\,\dd x$. Telle qu'écrite, l'intégrale intérieure
$\int\eu^{y^2}\dd y$ n'a pas de primitive élémentaire~: le
calcul est bloqué. Mais le domaine est le triangle $0 \leq
x \leq y \leq 1$, qui est élémentaire dans les deux directions~;
en échangeant l'\omterm{def:b2:structures:generated}{ordre},
\[
I = \int_0^1\!\!\int_0^y \eu^{y^2}\,\dd x\,\dd y
= \int_0^1 y\,\eu^{y^2}\,\dd y
= \Bigl[\tfrac12\eu^{y^2}\Bigr]_0^1 = \frac{\eu - 1}{2} .
\]
La variable intérieure $x$ n'apparaissait nulle part dans l'intégrande,
donc l'intégrer en premier a produit exactement le facteur $y$ qui
rend l'intégrale extérieure immédiate. Morale~: Fubini n'est pas
seulement une autorisation d'itérer --- c'est une autorisation de
\emph{choisir}, et le bon \omterm{def:b2:structures:generated}{ordre} peut transformer une intégrale
impossible en une ligne. Tracez toujours le domaine et lisez-en les
deux descriptions avant de commencer.
\end{example}

\begin{theorem}[Changement de variables]\label{thm:b2:multint:changevar}
Soit $\Phi \colon U' \to U$ un difféomorphisme $\mathcal{C}^1$
entre \omterm{def:b2:metric:topology}{ouverts} de $\R^2$, soit $K \subseteq U$ un domaine \omterm{def:b2:metric:compact}{compact}
découpé en morceaux élémentaires avec $K' = \Phi^{-1}(K)$, et
soit $f$ \omterm{def:b2:metric:continuity}{continue} sur $K$. Alors
\[
\iint_K f(x, y)\,\dd x\,\dd y
= \iint_{K'} f\bigl(\Phi(u, v)\bigr)\,
\abs{\det J_\Phi(u, v)}\,\dd u\,\dd v .
\]
\end{theorem}
\admitted

\begin{remark}
La preuve complète --- approcher $\Phi$ par sa \omterm{def:b2:diffcalc:differential}{différentielle} sur
une grille fine et contrôler les cellules du bord --- est longue mais
peu profonde~; elle est faite en entier dans la théorie de la mesure de
Licence~3, comme conséquence de la théorie de Lebesgue. L'heuristique
est l'image déjà utilisée pour l'\omterm{def:b2:multint:domain}{aire} d'une surface~: un petit carré
de côté $\dd u$ en $(u, v)$ est envoyé, au premier \omterm{def:b2:structures:generated}{ordre}, sur le
parallélogramme engendré par $\Phi_u\,\dd u$ et $\Phi_v\,\dd v$, dont
l'\omterm{def:b2:multint:domain}{aire} est $\abs{\det J_\Phi}\,\dd u\,\dd v$
(\cref{lem:b2:surfaces:lagrange}).
\end{remark}

\begin{remark}[Méthode~: choisir le changement de variables]
Trois réflexes couvrent la plupart des cas. \emph{Symétrie de
l'intégrande}~: $x^2 + y^2$ appelle les coordonnées polaires, une
structure de produit appelle à conserver les axes cartésiens.
\emph{Forme du bord}~: des bords $u(x,y) = c_1$, $v(x,y) = c_2$
réclament les coordonnées $(u, v)$ elles-mêmes, comme dans l'exemple
de la région hyperbolique ci-dessous --- le domaine devient un
rectangle, ce qui est toute la victoire. \emph{Structure linéaire}~:
les expressions en $x + y$ et $x - y$ invitent à la rotation de
$45$ degrés ou à une transvection (\cref{ex:b2:multint:affine}). Dans
tous les cas, trois cases à cocher avant d'intégrer~: l'application
est une bijection du nouveau domaine sur l'ancien~; son jacobien
est calculé \emph{dans le sens effectivement utilisé} (inverser à
la fin si c'est plus simple)~; et le jacobien intervient avec sa
valeur absolue.
\end{remark}

\begin{example}[Coordonnées polaires]\label{ex:b2:multint:polar}
$\Phi(\rho, \alpha) = (\rho\cos\alpha,\ \rho\sin\alpha)$ a
\[
J_\Phi = \begin{pmatrix}
\cos\alpha & -\rho\sin\alpha\\
\sin\alpha & \rho\cos\alpha
\end{pmatrix},
\qquad \det J_\Phi = \rho ,
\]
donc $\dd x\,\dd y = \rho\,\dd\rho\,\dd\alpha$. Pour le disque $D_R$
de rayon $R$~:
\[
\iint_{D_R} e^{-(x^2 + y^2)}\,\dd x\,\dd y
= \int_0^{2\pi}\!\!\int_0^R e^{-\rho^2}\rho\,\dd\rho\,\dd\alpha
= \pi\bigl(1 - e^{-R^2}\bigr)
\xrightarrow[R\to\infty]{} \pi .
\]
La comparaison avec le carré $[-R, R]^2$ (encadré entre les disques
$D_R$ et $D_{R\sqrt2}$, tous les intégrandes étant positifs) donne
$\bigl(\int_{-\infty}^\infty e^{-x^2}\dd x\bigr)^2 = \pi$~:
\[
\boxed{\ \int_{-\infty}^{+\infty} e^{-x^2}\,\dd x = \sqrt{\pi}\ }
\]
--- à nouveau l'intégrale de Gauss, cette fois par sa preuve la plus
célèbre (comparer avec le calcul à une variable du~%
\cref{ch:b2:integration}).
\end{example}

\begin{example}[Changements de variables affines]\label{ex:b2:multint:affine}
Pour une \omterm{def:b2:affine:subspace}{application affine} $\Phi(u, v) = M(u, v)^{\mathsf T} + C$ avec
$M$ inversible, le jacobien est la matrice constante $M$~:
les \omterm{def:b2:multint:domain}{aires} sont multipliées par le facteur constant $\abs{\det M}$
--- la promesse faite au~\cref{ch:b2:affine} est maintenant un théorème.
Deux usages immédiats. L'ellipse $\frac{x^2}{a^2} +
\frac{y^2}{b^2} \leq 1$ est l'image du disque unité par
$(u, v) \mapsto (au, bv)$, donc son \omterm{def:b2:multint:domain}{aire} est $ab \cdot \pi$ ---
sans calcul. Et pour l'intégrale de $f(x + y)$ sur le
carré $K = \intcc01^2$, la transvection $\Phi(u, v) = (u - v, v)$
(de \omterm{def:b2:linalg:det}{déterminant} $1$) la transforme en une intégrale de $f(u)$ sur un
parallélogramme, que Fubini découpe à $u$ constant~: avec $f =
\exp$,
\[
\iint_K \eu^{x+y}\,\dd x\,\dd y
= \Bigl(\int_0^1 \eu^x\,\dd x\Bigr)^2 = (\eu - 1)^2,
\]
comme le confirme la structure de produit. Choisir des coordonnées
adaptées à l'intégrande --- et non au domaine --- est l'autre
moitié du métier.
\end{example}


\begin{example}[Coordonnées adaptées à un domaine
curviligne]\label{ex:b2:multint:hyperbolicregion}
Soit $D$ la région du premier quadrant délimitée par les
hyperboles $xy = 1$ et $xy = 3$ et les droites $y = x$ et
$y = 3x$. Dans les coordonnées $u = xy$, $v = y/x$, le domaine
devient le carré $\intcc13 \times \intcc13$~; en inversant,
\[
x = \sqrt{u/v}, \qquad y = \sqrt{uv},
\qquad
\det J = x_uy_v - x_vy_u = \frac{1}{2v}
\]
(un calcul de deux lignes avec $x = u^{1/2}v^{-1/2}$, $y =
u^{1/2}v^{1/2}$). D'où
\[
\operatorname{Area}(D)
= \int_1^3\!\!\int_1^3\frac{\dd u\,\dd v}{2v}
= 2\cdot\frac{\ln 3}{2} = \ln 3 \approx 1.10 .
\]
Vouloir découper $D$ en coordonnées cartésiennes revient à le couper
en trois morceaux à bords hyperboliques et linéaires ---
faisable, sans joie, et propice aux erreurs. La morale reprend
l'\cref{ex:b2:multint:affine} dans toute sa force~: lisez les
équations du bord, et laissez-\emph{les} choisir les
coordonnées~; le jacobien convertit alors l'\omterm{def:b2:multint:domain}{aire} de la cellule du
maillage curviligne, exactement comme $\rho$ le faisait pour les
coordonnées polaires.
\end{example}

\begin{example}[Valeurs moyennes]\label{ex:b2:multint:average}
La \emph{valeur moyenne} de $f$ sur un domaine $D$ est
$\frac1{\operatorname{Area}(D)}\iint_Df$. Exemple~: la distance
moyenne au centre pour un point choisi \omterm{def:b2:funcseq:def}{uniformément} dans le
disque de rayon $R$ est
\[
\frac{1}{\pi R^2}\int_0^{2\pi}\!\!\int_0^R
\rho\cdot\rho\,\dd\rho\,\dd\alpha
= \frac{2\pi R^3/3}{\pi R^2} = \frac{2R}3 ,
\]
et non $R/2$~: l'\omterm{def:b2:multint:domain}{aire} uniforme met plus de masse aux grands rayons
(l'anneau au rayon $\rho$ a un poids proportionnel à $\rho$),
donc la moyenne se situe au-delà de la mi-distance. Obtenir ce
facteur correctement, c'est exactement le jacobien polaire à
l'œuvre, et la même pondération explique que le centre de gravité
$\bar z = 3R/8$ de la demi-boule, calculé plus loin dans le
chapitre, ne soit pas $R/2$.
\end{example}

\section{Le théorème de Green-Riemann}

\begin{theorem}[Green-Riemann]\label{thm:b2:multint:green}
Soit $K \subseteq \R^2$ un domaine \omterm{def:b2:metric:compact}{compact} élémentaire dans les
deux directions (ou une réunion finie de tels domaines recollés le
long de segments), de bord $\partial K$ une courbe \omterm{def:b2:multint:exact}{fermée}
$\mathcal{C}^1$ par morceaux orientée dans le \emph{sens direct} (le
domaine reste à gauche). Pour $P, Q$ de classe $\mathcal{C}^1$ sur un
voisinage de $K$~:
\[
\oint_{\partial K} P\,\dd x + Q\,\dd y
= \iint_K \Bigl(\frac{\partial Q}{\partial x}
- \frac{\partial P}{\partial y}\Bigr)\,\dd x\,\dd y .
\]
\end{theorem}

\begin{proof}
D'abord, les deux membres sont additifs lorsqu'on coupe $K$ le long
d'un segment en deux morceaux $K_1, K_2$~: les intégrales doubles
s'ajoutent par additivité de $\iint$~; quant aux intégrales de bord,
les bords orientés dans le sens direct de $K_1$ et $K_2$ parcourent
chacun la coupe intérieure une fois, en sens \emph{opposés}, de sorte
que dans la somme
\[
\oint_{\partial K_1} + \oint_{\partial K_2}
= \oint_{\partial K} + (\text{coupe, dans les deux sens})
= \oint_{\partial K},
\]
les deux passages le long de la coupe s'annulent
(\cref{prop:b2:multint:lineinv}) et seul le bord extérieur subsiste.
En itérant un nombre fini de coupes, il suffit de traiter un domaine
élémentaire. Prouvons $\oint P\,\dd x = -\iint_K P_y$ sur un domaine
$D = \{a \leq x \leq b,\ \varphi_1(x) \leq y \leq \varphi_2(x)\}$
élémentaire selon $y$~; l'identité $\oint Q\,\dd y = \iint_K Q_x$
est \omterm{def:b2:quadratic:adjoint}{symétrique} (élémentaire selon $x$), et le théorème en est la somme.

Calculons l'intégrale double par Fubini et le théorème fondamental à
une variable~:
\[
\iint_D \frac{\partial P}{\partial y}\,\dd x\,\dd y
= \int_a^b \bigl(P(x, \varphi_2(x)) - P(x, \varphi_1(x))\bigr)
\,\dd x .
\]
Maintenant le bord de $D$, dans le sens direct, se compose de~: le
graphe inférieur $y = \varphi_1(x)$ parcouru de gauche à droite, le
segment vertical droit $x = b$ (vers le haut), le graphe supérieur
$y = \varphi_2(x)$ parcouru \emph{de droite à gauche}, le segment
vertical gauche $x = a$ (vers le bas). Le long des segments verticaux,
$x$ est constant, donc ils contribuent $0$ à $\oint P\,\dd x$~; les
graphes, paramétrés par $x$, donnent
\[
\oint_{\partial D} P\,\dd x
= \int_a^b P(x, \varphi_1(x))\,\dd x
- \int_a^b P(x, \varphi_2(x))\,\dd x
= -\iint_D \frac{\partial P}{\partial y}\,\dd x\,\dd y .
\qedhere
\]
\end{proof}

\begin{corollary}[Aire par le bord]\label{cor:b2:multint:area}
Sous les hypothèses du~\cref{thm:b2:multint:green},
\[
\operatorname{Area}(K)
= \oint_{\partial K} x\,\dd y
= -\oint_{\partial K} y\,\dd x
= \frac12\oint_{\partial K} x\,\dd y - y\,\dd x .
\]
\end{corollary}

\begin{proof}
Appliquer Green-Riemann à $(P, Q) = (0, x)$, $(-y, 0)$ et
$\frac12(-y, x)$~: à chaque fois $Q_x - P_y = 1$.
\end{proof}

\begin{remark}[Choisir parmi les trois formules d'aire]
Les trois formules de bord sont égales, non interchangeables en
pratique. Utilisez $\oint x\,\dd y$ quand le paramétrage rend
$\dd y$ simple (graphes au-dessus de l'axe des $y$), $-\oint y\,\dd x$
symétriquement, et la demi-somme \omterm{def:b2:quadratic:adjoint}{symétrique} quand le
paramétrage traite $x$ et $y$ de façon équilibrée --- pour l'ellipse
elle a produit un intégrande constant, sans aucune linéarisation
trigonométrique. Sur les bords polygonaux, la demi-somme
devient la formule du lacet de l'\cref{exo:b2:multint:12},
l'algorithme des arpenteurs. Et lorsque le bord est parcouru
dans le sens horaire par le paramétrage donné, les trois formules
renvoient \emph{moins} l'\omterm{def:b2:multint:domain}{aire}~: un résultat négatif n'est pas une
erreur de calcul mais un rapport d'orientation --- changez le
signe, ou le paramétrage.
\end{remark}

\begin{example}[Aire de l'ellipse]\label{ex:b2:multint:ellipsearea}
Pour $x = a\cos t$, $y = b\sin t$, $t \in [0, 2\pi]$~:
\[
\operatorname{Area}
= \frac12\int_0^{2\pi}\bigl(a\cos t \cdot b\cos t
- b\sin t\cdot(-a\sin t)\bigr)\,\dd t
= \frac{ab}{2}\int_0^{2\pi}\dd t = \pi ab .
\]
\end{example}

\begin{example}[Green-Riemann comme
vérification croisée]\label{ex:b2:multint:greencheck}
Prenons $P = -y^3$, $Q = x^3$ sur le disque unité fermé $D$.
Côté bord, avec $\gamma(t) = (\cos t, \sin t)$~:
\[
\oint_{\partial D}P\,\dd x + Q\,\dd y
= \int_0^{2\pi}\bigl(\sin^4 t + \cos^4 t\bigr)\dd t
= 2\pi\cdot\Bigl(\frac38 + \frac38\Bigr) = \frac{3\pi}2 ,
\]
par linéarisation ($\sin^4 + \cos^4 = \tfrac34 +
\tfrac14\cos4t$). Côté intérieur~:
\[
\iint_D(Q_x - P_y)\,\dd x\,\dd y
= \iint_D 3(x^2 + y^2)\,\dd x\,\dd y
= 3\int_0^{2\pi}\!\!\int_0^1\rho^3\,\dd\rho\,\dd\alpha
= \frac{3\pi}2 .
\]
Même nombre, deux calculs très différents --- et c'est là
l'usage pratique~: le membre le plus facile de l'identité de Green
devient le calcul, l'autre une vérification. Pour les
circulations de champs polynomiaux le long de courbes \omterm{def:b2:multint:exact}{fermées}, l'%
intégrale double est presque toujours le membre facile.
\end{example}

\begin{remark}
Green-Riemann explique l'\cref{ex:b2:multint:angleform}~: pour une
\omterm{def:b2:multint:exact}{forme fermée} ($Q_x = P_y$), l'intégrale sur le bord de tout domaine
contenu dans $U$ est nulle. La forme angulaire n'est pas \omterm{def:b2:multint:exact}{exacte}
uniquement parce que le point retiré à l'origine empêche le disque
bordé par le cercle unité de se trouver dans $U$ --- les \omterm{def:b2:multint:lineint}{intégrales
curvilignes} des \omterm{def:b2:multint:exact}{formes fermées} détectent les trous du domaine.
(Poussée plus loin, cette observation devient la cohomologie de de
Rham.)
\end{remark}

\section{Intégrales triples}

La théorie s'étend à trois variables sans idée nouvelle~: Fubini
ramène $\iiint$ à trois intégrales à une variable (soit par
\emph{tranches}~: $\iiint_K f = \int\bigl(\iint_{K_z}
f\bigr)\dd z$ sur les tranches horizontales $K_z$, soit par
\emph{empilement}~: en intégrant d'abord en $z$ le long de bâtonnets
verticaux), et la formule de changement de variables vaut avec le
jacobien $3 \times 3$.

\begin{example}[Coordonnées cylindriques et sphériques]\label{ex:b2:multint:spherical}
\emph{Cylindriques} $(x, y, z) = (\rho\cos\alpha, \rho\sin\alpha,
z)$~: $\dd x\,\dd y\,\dd z = \rho\,\dd\rho\,\dd\alpha\,\dd z$.
\emph{Sphériques} $(x, y, z) = (r\cos\theta\cos\varphi,\
r\sin\theta\cos\varphi,\ r\sin\varphi)$ ($\theta$ longitude,
$\varphi \in [-\frac\pi2, \frac\pi2]$ latitude)~: en développant le $3
\times 3$ \omterm{def:b2:linalg:det}{déterminant} selon la dernière ligne,
\[
\det J = r^2\cos\varphi ,
\qquad
\dd x\,\dd y\,\dd z
= r^2\cos\varphi\;\dd r\,\dd\theta\,\dd\varphi .
\]
\omterm{pb:b2:multint:1}{Volume de la boule} de rayon $R$~:
\[
V = \int_0^R\!\!\int_0^{2\pi}\!\!\int_{-\pi/2}^{\pi/2}
r^2\cos\varphi\;\dd\varphi\,\dd\theta\,\dd r
= \frac{R^3}{3}\cdot 2\pi \cdot 2
= \boxed{\frac43\pi R^3} ,
\]
acquittant enfin la formule admise dans les chapitres de volume des
livres précédents.
\end{example}

\begin{example}[Le tétraèdre, deux fois]
\label{ex:b2:multint:tetra}
Le volume de $T = \{x, y, z \geq 0,\ x + y + z \leq 1\}$, par
empilement~: pour $(x, y)$ fixé dans le triangle $x + y \leq 1$,
$z$ parcourt $\intcc0{1 - x - y}$, donc
\[
V = \int_0^1\!\!\int_0^{1-x}(1 - x - y)\,\dd y\,\dd x
= \int_0^1\frac{(1 - x)^2}{2}\,\dd x = \frac16 .
\]
Par tranches~: la section à hauteur $z$ est le triangle $\{x, y
\geq 0,\ x + y \leq 1 - z\}$, d'\omterm{def:b2:multint:domain}{aire} $\frac{(1-z)^2}2$, et
$V = \int_0^1\frac{(1-z)^2}2\,\dd z = \frac16$ à nouveau --- les
deux calculs sont les mêmes intégrales dans un \omterm{def:b2:structures:generated}{ordre} différent,
ce qui est tout ce qu'affirme Fubini. La valeur $\frac16 =
\frac13\cdot\frac12\cdot1$ est la formule du cône
(\cref{ex:b2:multint:cone}) à base triangulaire, et la
version en dimension $n$, $1/n!$, se démontre par exactement ce
découpage en tranches dans le problème du week-end.
\end{example}

\begin{example}[Centre de gravité d'une demi-boule]\label{ex:b2:multint:centroid}
Pour la demi-boule supérieure $H$ de rayon $R$ ($z \geq 0$), la
hauteur du centre de gravité est $\bar z = \frac1{V}\iiint_H z$, avec
$V = \frac23\pi R^3$. En coordonnées sphériques ($z =
r\sin\varphi$, $\varphi \in \intcc0{\pi/2}$)~:
\[
\iiint_H z
= \int_0^R r^3\,\dd r\int_0^{2\pi}\dd\theta
\int_0^{\pi/2}\sin\varphi\cos\varphi\,\dd\varphi
= \frac{R^4}4\cdot2\pi\cdot\frac12 = \frac{\pi R^4}4 ,
\]
donc
\[
\bar z = \frac{\pi R^4/4}{2\pi R^3/3} = \frac{3R}8 :
\]
le point d'équilibre d'un hémisphère plein se situe à trois huitièmes
du rayon au-dessus de la face plate --- en dessous de la mi-hauteur
$R/2$, comme il se doit, puisque le solide est plus épais près de la
base. Tout calcul de centre de gravité a cette forme~: une intégrale
de moment, un volume, un rapport, et un contrôle de vraisemblance
face à la géométrie.
\end{example}

\begin{example}[Volume par tranches~: le cône]\label{ex:b2:multint:cone}
Un cône d'\omterm{def:b2:multint:domain}{aire} de base $A$ et de hauteur $h$ (sommet en haut, base à $z = 0$)~:
la tranche à hauteur $z$ est la base dilatée du facteur $(1 -
z/h)$, d'\omterm{def:b2:multint:domain}{aire} $A(1 - z/h)^2$. D'où
\[
V = \int_0^h A\Bigl(1 - \frac zh\Bigr)^2\dd z = \frac{Ah}{3} :
\]
le un tiers des formules du lycée, valable pour \emph{n'importe quelle}
forme de base --- le découpage en tranches la transforme en l'intégrale
d'un carré.
\end{example}

\begin{example}[Seuils d'intégrabilité dans le
plan]\label{ex:b2:multint:threshold}
Pour quels $\alpha > 0$ l'intégrale $\iint_{D}\rho^{-\alpha}\,\dd
x\,\dd y$ converge-t-elle sur le disque unité privé de l'origine $D$
(limite sur les anneaux $\varepsilon \leq \rho \leq 1$)~? En
coordonnées polaires,
\[
\int_0^{2\pi}\!\!\int_\varepsilon^1\rho^{-\alpha}\,
\rho\,\dd\rho\,\dd\alpha
= 2\pi\int_\varepsilon^1\rho^{1-\alpha}\,\dd\rho ,
\]
qui \omterm{def:b2:integration:improper}{converge} lorsque $\varepsilon \to 0$ si et seulement si
$1 - \alpha > -1$, c.-à-d.\ $\alpha < 2$~: en dimension $2$ l'exposant
critique de singularité est la dimension elle-même, le $\rho$
supplémentaire venu du jacobien adoucissant la singularité d'une
puissance. (De même $\alpha < 3$ pour une singularité ponctuelle dans
l'espace, via $r^2$.) Une comptabilité radiale de ce type est la
façon dont l'intégrabilité se décide d'un coup d'œil dans le cadre
de Lebesgue de Licence~3 --- et c'est la raison pour laquelle
$\iiint 1/r$ convergeait sans effort à l'\cref{exo:b2:multint:7}.
\end{example}

\begin{remark}[Pièges courants]
(i) \emph{Orientation}~: une \omterm{def:b2:multint:lineint}{intégrale curviligne} change de signe avec
le sens de parcours, et Green-Riemann exige le bord dans le sens
direct (domaine à gauche)~; pour un domaine avec un trou, le bord
intérieur est parcouru dans le \emph{sens horaire}. (ii)
\emph{Le jacobien intervient avec une valeur absolue}~: le changement
de variables ne produit jamais d'\omterm{def:b2:multint:domain}{aire} négative, et oublier
$\abs{\det}$ change généralement les signes précisément quand
l'application renverse l'orientation. (iii) \emph{Le facteur polaire
$\rho$}~: $\dd x\,\dd y = \rho\,\dd\rho\,\dd\alpha$, et non
$\dd\rho\,\dd\alpha$ --- l'erreur la plus fréquente de tout le
chapitre~; l'analyse dimensionnelle la détecte, puisque
$\dd\rho\,\dd\alpha$ a la dimension d'une \omterm{def:b2:curves:length}{longueur}, non d'une
\omterm{def:b2:multint:domain}{aire}. (iv) \emph{Intégrales doubles impropres}~: les limites sur des
disques croissants et sur des carrés croissants coïncident ici parce
que les intégrandes sont positifs (encadrement)~; pour des intégrandes
changeant de signe, la limite peut dépendre de l'exhaustion, et aucune
affirmation n'est faite sans convergence absolue. (v)
\emph{Domaines contre intégrandes}~: un intégrande produit sur un
domaine non produit ne factorise \emph{pas} l'intégrale ---
la factorisation exige les deux, comme dans le carré de
l'\cref{ex:b2:multint:polar}.
\end{remark}

\begin{remark}[Perspectives au sein de ce volume]
L'intégrale de Gauss calculée ici est discrètement partout dans
les chapitres de probabilités~: la constante $\sqrt\pi$ à l'intérieur
de la formule de Stirling (\cref{thm:b2:comparison:stirling}) est
l'intégrale de ce chapitre, et via Stirling elle fixe l'asymptotique
en $1/\sqrt{\pi n}$ des probabilités de retour de la marche
aléatoire au~\cref{ch:b2:proba}. Les \omterm{pb:b2:multint:1}{intégrales de Wallis} du
problème du week-end y réapparaissent aussi, gouvernant les mêmes
estimations du coefficient binomial central. Dans l'autre sens, les
éléments d'\omterm{def:b2:multint:domain}{aire} et de volume de ce chapitre complètent la géométrie
du~\cref{ch:b2:surfaces}, et la formule de Green recalcule les
\omterm{def:b2:multint:domain}{aires} des \omterm{pb:b2:curves:1}{enveloppes} du~\cref{ch:b2:curves} (l'\omterm{pb:b2:curves:1}{astroïde}, à
l'\cref{exo:b2:multint:5}). Un chapitre, trois services~:
la mesure pour la géométrie, les constantes pour les probabilités, et
la discipline du changement de variables utilisée par les deux.
\end{remark}

\section{Exercices}

\begin{exercise}[$\star$]\label{exo:b2:multint:1}
Calculer $\int_\gamma y^2\,\dd x + x\,\dd y$ le long de~: (a) le segment
de $(0,0)$ à $(1,1)$~; (b) l'arc de parabole $y = x^2$ de
$(0,0)$ à $(1,1)$. La forme est-elle \omterm{def:b2:multint:exact}{exacte}~?
\end{exercise}

\begin{exercise}[$\star$]\label{exo:b2:multint:2}
Montrer que $\omega = (2xy + y^3)\,\dd x + (x^2 + 3xy^2 + 1)\,\dd y$
est \omterm{def:b2:multint:exact}{fermée} sur $\R^2$, trouver un \omterm{def:b2:multint:exact}{potentiel}, et calculer
$\int_\gamma\omega$ le long de tout arc de $(0, 0)$ à $(1, 2)$.
\end{exercise}

\begin{exercise}[$\star$]\label{exo:b2:multint:3}
Calculer $\iint_D (x + y)\,\dd x\,\dd y$ où $D$ est le domaine
délimité par $y = x^2$ et $y = x$ ($0 \leq x \leq 1$), dans les deux
\omterm{def:b2:structures:generated}{ordres} d'intégration.
\end{exercise}

\begin{exercise}[$\star\star$]\label{exo:b2:multint:4}
En utilisant les coordonnées polaires, calculer $\iint_D \frac{\dd x\,\dd y}{(1 +
x^2 + y^2)^2}$ sur le plan entier (comme limite sur des disques), et
$\iint_{D'} xy\,\dd x\,\dd y$ sur le quart de disque $D' = \{x, y
\geq 0,\ x^2 + y^2 \leq 1\}$.
\end{exercise}

\begin{exercise}[$\star\star$]\label{exo:b2:multint:5}
Calculer l'\omterm{def:b2:multint:domain}{aire} délimitée par l'\omterm{pb:b2:curves:1}{astroïde} $x = \cos^3 t$, $y =
\sin^3 t$, $t \in [0, 2\pi]$, en utilisant
le~\cref{cor:b2:multint:area}. \emph{(Linéariser $\sin^2 t\cos^2 t$.)}
\end{exercise}

\begin{exercise}[$\star\star$]\label{exo:b2:multint:6}
Calculer le volume du solide délimité par le bas par le paraboloïde $z
= x^2 + y^2$ et par le haut par le plan $z = 1$, par les deux méthodes~:
empilement (intégrer $1 - x^2 - y^2$ sur le disque unité, coordonnées
polaires) et tranches (les tranches horizontales sont des disques de
rayon $\sqrt z$).
\end{exercise}

\begin{exercise}[$\star\star$]\label{exo:b2:multint:7}
(Attraction gravitationnelle d'une boule --- théorème de Newton, cas
particulier) Montrer que le volume de la coquille sphérique $a \leq r \leq
b$ est $\frac43\pi(b^3 - a^3)$ et calculer $\iiint_{B}
\frac{\dd x\,\dd y\,\dd z}{r}$ sur la boule $B$ de rayon $R$
($r$ la distance à l'origine). \emph{(Coordonnées
sphériques.)}
\end{exercise}

\begin{exercise}[$\star\star\star$]\label{exo:b2:multint:8}
(Lemme de Poincaré, cas étoilé) Soit $U$ étoilé par rapport
à $0$ (c.-à-d.\ $M \in U \Rightarrow [0, M] \subseteq U$)
et $\omega = P\dd x + Q\dd y$ une forme $\mathcal{C}^1$ \omterm{def:b2:multint:exact}{fermée} sur
$U$. On pose
\[
f(x, y) = \int_0^1 \bigl(x\,P(tx, ty) + y\,Q(tx, ty)\bigr)\dd t .
\]
En utilisant la dérivation sous le signe intégral
(\cref{ch:b2:integration}) et $P_y = Q_x$, montrer que $f_x = P$
et $f_y = Q$~: toute \omterm{def:b2:multint:exact}{forme fermée} sur un \omterm{def:b2:metric:topology}{ouvert} étoilé est
\omterm{def:b2:multint:exact}{exacte}.
\end{exercise}

\begin{exercise}[$\star\star\star$]\label{exo:b2:multint:9}
(Intégrale de Dirichlet par intégration double) Justifier et exploiter
\[
\int_0^\infty\!\!\int_0^\infty e^{-xy}\sin x\;\dd y\,\dd x
\quad\text{vs}\quad
\int_0^\infty\!\!\int_0^\infty e^{-xy}\sin x\;\dd x\,\dd y
\]
sur $[0, A] \times [0, \infty)$~: montrer $\int_0^A \frac{\sin
x}{x}\dd x = \frac\pi2 - \int_0^\infty
e^{-Ay}\frac{y\sin A + \cos A}{1 + y^2}\dd y$ et retrouver
$\int_0^\infty \frac{\sin x}{x}\,\dd x = \frac\pi2$, en comparant
avec la preuve par \omterm{thm:b2:integration:continuity}{intégrale à paramètre} du~%
\cref{ch:b2:integration}.
\end{exercise}

\begin{exercise}[$\star\star\star$]\label{exo:b2:multint:10}
(Inégalité isopérimétrique via Wirtinger) Soit $\gamma$ une courbe
\omterm{def:b2:multint:exact}{fermée} simple $\mathcal{C}^1$ de \omterm{def:b2:curves:length}{longueur} $2\pi$, paramétrée par
l'\omterm{def:b2:curves:length}{abscisse curviligne} sur $[0, 2\pi]$, délimitant une \omterm{def:b2:multint:domain}{aire} $A$. En
utilisant le~\cref{cor:b2:multint:area}, Parseval et l'inégalité de
Wirtinger (exercices du~\cref{ch:b2:fourier}), prouver $A \leq \pi$, avec
égalité pour le cercle. \emph{(Normaliser $\int_0^{2\pi} x(s)\dd s
= 0$~; écrire $2A = \oint x\,\dd y - y\,\dd x$ et majorer $2A \leq
\int (x^2 + y'^2)$ soigneusement via $2A = \int_0^{2\pi}(xy' - yx')\dd
s$ et $x^2 + y'^2 \geq 2xy'$.)}
\end{exercise}

\begin{exercise}[$\star\star$]\label{exo:b2:multint:11}
(Moments de la boule) Pour la boule $B$ de rayon $R$ dans
$\R^3$, calculer $\iiint_B z^2\,\dd x\,\dd y\,\dd z$ en
coordonnées sphériques, et en déduire $\iiint_B (x^2 + y^2 +
z^2)\,\dd x\,\dd y\,\dd z$ par symétrie. Vérifier ce
dernier par le calcul en coquilles $\int_0^R r^2\cdot4\pi
r^2\,\dd r$.
\end{exercise}

\begin{exercise}[$\star\star$]\label{exo:b2:multint:12}
(Formule du lacet) Soit $K$ un polygone de sommets $(x_1,
y_1), \dots, (x_m, y_m)$ dans le sens direct (indices
modulo $m$). Déduire du~\cref{cor:b2:multint:area} que
\[
\operatorname{Area}(K)
= \frac12\sum_{i=1}^m
\bigl(x_iy_{i+1} - x_{i+1}y_i\bigr) ,
\]
et vérifier la formule sur le triangle $(0,0)$, $(1,0)$,
$(0,1)$.
\end{exercise}

\section{Problème~: le volume de la boule en dimension $n$}

\begin{problem}[{Problème du week-end --- $V_n =
\pi^{n/2}/\Gamma(\frac n2 + 1)$, et l'étrangeté des grandes
dimensions}]\label{pb:b2:multint:1}
Le disque a pour \omterm{def:b2:multint:domain}{aire} $\pi$, la boule pour volume $\frac43\pi$ --- et
ensuite~? Ce problème calcule le \omterm{pb:b2:multint:1}{volume de la boule} unité de
$\R^n$ pour tout $n$, deux fois (par une récurrence par tranches
guidée par les \omterm{pb:b2:multint:1}{intégrales de Wallis}, puis via la fonction $\Gamma$ et
l'intégrale de Gauss de l'\cref{ex:b2:multint:polar}), puis en lit
la géométrie~: les volumes culminent en dimension cinq et se
précipitent vers zéro, et presque toute une boule de grande dimension
se cache dans une mince coquille près de son bord.\index{volume de la
boule}\index{intégrales de Wallis} Pour une fonction \omterm{def:b2:metric:continuity}{continue} sur une
boule de $\R^n$, l'intégrale s'entend comme l'intégrale itérée $n$~fois
(en tranchant une coordonnée à la fois, comme dans le chapitre pour $n
\leq 3$)~; on note $B_n(R)$ la boule \omterm{def:b2:multint:exact}{fermée} de rayon $R$
centrée en $0$, $v_n(R)$ son volume, et $V_n =
v_n(1)$, avec $V_0 = 1$ par convention.

\medskip
\noindent\textbf{Partie I --- La récurrence par tranches.}
\begin{enumerate}
  \item En substituant $x_i = Ru_i$ dans chacune des $n$ intégrales
        itérées, montrer que $v_n(R) = V_nR^n$.
  \item En tranchant $B_n(1)$ selon sa dernière coordonnée, montrer
        que
        \[
        V_n = V_{n-1}\int_{-1}^{1}(1 -
        t^2)^{\frac{n-1}2}\,\dd t .
        \]
  \item Avec $t = \sin\theta$, identifier l'intégrale à une
        intégrale de Wallis~: $\int_{-1}^1(1 -
        t^2)^{\frac{n-1}2}\dd t = 2W_n$, où $W_n =
        \int_0^{\pi/2}\cos^n\theta\,\dd\theta =
        \int_0^{\pi/2}\sin^n\theta\,\dd\theta$.
  \item Prouver les deux identités de Wallis (intégrer par parties~;
        puis télescoper $nW_nW_{n-1}$)~:
        \[
        W_n = \frac{n-1}nW_{n-2} \quad (n \geq 2),
        \qquad
        W_nW_{n-1} = \frac{\pi}{2n} \quad (n \geq 1).
        \]
\end{enumerate}

\medskip
\noindent\textbf{Partie II --- La récurrence résolue.}
\begin{enumerate}[resume]
  \item Combiner les questions 2 à 4 en la récurrence à deux pas
        \[
        V_n = \frac{2\pi}{n}\,V_{n-2}
        \qquad (n \geq 2).
        \]
  \item En déduire les expressions explicites, pour $k \geq 0$~:
        \[
        V_{2k} = \frac{\pi^k}{k!},
        \qquad
        V_{2k+1} = \frac{2^{k+1}\pi^k}{1\cdot3\cdot5\cdots
        (2k+1)} .
        \]
  \item Tabuler $V_1, \dots, V_7$ numériquement. En utilisant le
        rapport $V_n/V_{n-2} = 2\pi/n$ et les valeurs de $2W_5$
        et $2W_6$, prouver que la suite $(V_n)$ croît
        jusqu'à son maximum $V_5 = \frac{8\pi^2}{15} \approx
        5.26$ puis décroît à partir de là.
  \item Montrer que $V_n \to 0$ plus vite que toute suite géométrique,
        et que $\sum_{n\geq1} V_n$ \omterm{def:b2:integration:improper}{converge}~: toutes
        les boules unité ensemble ont un volume total fini.
  \item Prouver l'identité génératrice
        \[
        \sum_{k\geq0} V_{2k}\,x^{2k} = \eu^{\pi x^2}
        \qquad (x \in \R),
        \]
        et en déduire $\sum_{k\geq0}V_{2k} = \eu^\pi \approx
        23.14$.
\end{enumerate}

\medskip
\noindent\textbf{Partie III --- Deuxième voie~: $\Gamma$ et
l'intégrale de Gauss.}
\begin{enumerate}[resume]
  \item Montrer par Fubini (l'intégrande est un produit) que
        \[
        I_n = \int_{\R^n}\eu^{-\norm
        x^2}\dd x
        = \Bigl(\int_{-\infty}^{+\infty}
        \eu^{-t^2}\dd t\Bigr)^{\!n} = \pi^{n/2},
        \]
        l'intégrale de Gauss en dimension $n$, entendue comme une
        limite sur les cubes $\intcc{-R}{R}^n$.
  \item Rappeler $\Gamma(s) = \int_0^\infty t^{s-1}\eu^{-t}\dd
        t$ (\cref{def:b2:integration:gamma}). À partir de
        $\Gamma(s+1) = s\,\Gamma(s)$
        (\cref{thm:b2:integration:gammaprops}) et
        $\Gamma(\tfrac12) = \sqrt\pi$ (substituer $t = u^2$
        et invoquer l'intégrale de Gauss), calculer
        \[
        \Gamma(k + 1) = k!,
        \qquad
        \Gamma\Bigl(k + \frac32\Bigr) =
        \frac{1\cdot3\cdots(2k+1)}{2^{k+1}}\,\sqrt\pi .
        \]
  \item Prouver, par récurrence via la récurrence de la question
        5, la formule unique
        \[
        V_n = \frac{\pi^{n/2}}{\Gamma\bigl(\frac n2 +
        1\bigr)} \qquad (n \geq 1),
        \]
        et vérifier qu'elle reproduit les deux expressions explicites
        de la question 6.
  \item Montrer $\int_0^\infty \eu^{-r^2}r^{n-1}\dd r =
        \tfrac12\Gamma\bigl(\tfrac n2\bigr)$ et en déduire
        l'identité
        \[
        I_n = n\,V_n\int_0^\infty \eu^{-r^2}\,r^{n-1}\,\dd r.
        \]
        L'interpréter~: la masse gaussienne de $\R^n$ se
        collecte le long de coquilles sphériques dont
        l'«~\omterm{def:b2:multint:domain}{aire} $(n-1)$-dimensionnelle~» au rayon $r$ vaut
        $nV_nr^{n-1}$ --- les deux membres sont désormais prouvés
        indépendamment, donc l'interprétation ne coûte rien.
  \item Poser $s_{n-1} = nV_n$ (l'\omterm{def:b2:multint:domain}{aire} de la sphère unité
        $S^{n-1}$, en cohérence avec $v_n(R) = \int_0^R
        s_{n-1}r^{n-1}\dd r$). Tabuler $s_0, \dots, s_3$ et
        vérifier $s_1 = 2\pi$, $s_2 = 4\pi$, $s_3 = 2\pi^2$.
\end{enumerate}

\medskip
\noindent\textbf{Partie IV --- Les grandes dimensions sont étranges.}
\begin{enumerate}[resume]
  \item À partir de la formule de Stirling
        (\cref{thm:b2:comparison:stirling}) appliquée à $k!$,
        montrer pour $n = 2k$ pair~:
        \[
        V_n \sim \frac{1}{\sqrt{\pi n}}
        \Bigl(\frac{2\pi\eu}{n}\Bigr)^{n/2}
        \qquad (n \to \infty, \ n \text{ pair}),
        \]
        et expliquer pourquoi la même borne de décroissance
        super-géométrique s'étend à $n$ impair via la récurrence.
  \item La boule unité est contenue dans le cube $\intcc{-1}1^n$ de
        volume $2^n$. Calculer le taux de remplissage $V_n/2^n$ pour
        $n = 2, 3, 10$, et montrer qu'il tend vers $0$~: en grande
        dimension, l'essentiel du cube se trouve dans ses coins.
  \item Montrer que la fraction de $v_n(1)$ située à
        distance au plus $\varepsilon$ de la sphère de bord est $1 -
        (1 - \varepsilon)^n \to 1$~; numériquement, quelle
        fraction d'une boule en dimension $100$ se trouve dans la
        coquille extérieure d'épaisseur $1\%$~?
  \item Prouver l'asymptotique de Wallis $W_n \sim
        \sqrt{\dfrac{\pi}{2n}}$ \emph{(monotonie de
        $(W_n)$, le rapport $W_n/W_{n-2} \to 1$, et $W_nW_{n-1}
        = \frac\pi{2n}$)}, et la minoration $W_n \geq
        \sqrt{\dfrac{\pi}{2(n+1)}}$ pour tout $n$.
  \item (Concentration sur une dalle) La fraction de la boule
        unité dont la première coordonnée dépasse $\delta$ est
        $\int_\delta^1(1 - x^2)^{\frac{n-1}2}\dd x \,\big/\,
        (2W_n)$. En utilisant $1 - u \leq \eu^{-u}$ et la
        majoration de queue $\int_\delta^\infty \eu^{-a x^2}\dd x \leq
        \frac{\eu^{-a\delta^2}}{2a\delta}$, montrer que cette
        fraction est au plus
        \[
        \frac{\eu^{-(n-1)\delta^2/2}}{(n-1)\,\delta}
        \Big/ \sqrt{\frac{2\pi}{n+1}}
        \]
        et conclure~: pour $\delta = s/\sqrt{n-1}$, tout sauf une
        fraction $O(\eu^{-s^2/2}/s)$ de la boule se trouve dans la
        dalle $\abs{x_1} \leq s/\sqrt{n-1}$. Une boule de grande
        dimension est, statistiquement, une mince crêpe dans toutes
        les directions à la fois.
  \item Assembler les questions 16 à 19 en un paragraphe~: où
        se situe le volume de $B_n(1)$ (près de la sphère de bord,
        tout en restant dans des dalles $O(1/\sqrt n)$ de tout
        hyperplan passant par le centre), et pourquoi ces deux
        affirmations ne se contredisent pas.
\end{enumerate}

\medskip
\noindent\textbf{Partie V --- Autres corps, et synthèse.}
\begin{enumerate}[resume]
  \item (Simplexe) Soit $\Delta_n = \{x \in \R^n : x_i \geq 0,\
        \sum x_i \leq 1\}$. Montrer par tranches et récurrence que
        $\operatorname{vol}(\Delta_n) = \frac1{n!}$.
  \item (Polytope croisé) En déduire que $C_n = \{x :
        \sum\abs{x_i} \leq 1\}$ a pour volume $\frac{2^n}{n!}$,
        et vérifier l'encadrement $C_n \subseteq B_n(1)
        \subseteq \intcc{-1}1^n$ au niveau des volumes~:
        $\frac{2^n}{n!} \leq V_n \leq 2^n$.
  \item Calculer $V_4$ d'une troisième façon~: trancher $\R^4 = \R^2 \times
        \R^2$, intégrer l'\omterm{def:b2:multint:domain}{aire} du disque en $(z, w)$ sur
        le disque en $(x, y)$ en coordonnées polaires, et retrouver
        $V_4 = \frac{\pi^2}2$.
  \item (Monte-Carlo en difficulté) Un point est tiré \omterm{def:b2:funcseq:def}{uniformément}
        dans le cube $\intcc{-1}1^{20}$. Montrer que la
        probabilité qu'il tombe dans la boule inscrite est
        $V_{20}/2^{20} \approx 2.5\cdot10^{-8}$, de sorte qu'environ
        quarante millions de tirages sont nécessaires avant d'espérer
        le premier succès~: estimer $V_n$ par échantillonnage par rejet
        s'effondre en grande dimension (le fléau de la
        dimension).
  \item Synthèse. Deux dérivations indépendantes se sont rejointes
        en $V_n = \pi^{n/2}/\Gamma(\frac n2 + 1)$~: lister quel
        théorème de ce chapitre chacune a utilisé (Fubini, changement
        de variables, l'intégrale de Gauss polaire), et quels
        ingrédients à une variable (Wallis, $\Gamma$, Stirling).
        Où le volume de Licence~3 refait-il ce calcul
        avec la théorie de Lebesgue, et qu'y ajoute-t-il~?
\end{enumerate}
\end{problem}
